\chapter{Supplementary Experimental Details}
\label{app:experimental-details}

This appendix records the experimental details that are compressed in
Chapter~\ref{chap:experiment-analysis}. The goal is to make the benchmark
construction, model configuration, training setup, and baseline adaptations
explicit enough for the results to be audited and reproduced.

\section{Training and Evaluation Split}
\label{app:training-evaluation-split}

The GRPO training data is constructed from three sources. LoCoMo and
LongMemEval provide public multi-session dialogue-memory questions with gold
answers and evidence-session annotations. The diagnostic \bench{} benchmark
contributes temporal-bridge records whose final answer evidence is not directly
addressable from the original query.

We start from 1,986 LoCoMo questions, 500 LongMemEval questions, and 500 raw
\bench{} candidates. For LoCoMo, category-5 adversarial questions are excluded.
For LongMemEval, single-session preference questions are excluded. For
\bench{}, the generation pipeline removes invalid attempts, duplicates, split
leakage, incomplete metadata, temporal inconsistencies, and non-unique answers.
This produces 1,536 LoCoMo questions, 469 LongMemEval questions, and 374
\bench{} records in the filtered candidate pool.

The final GRPO training set contains 500 examples. To avoid conversation-level
leakage, LoCoMo is partitioned by dialogue rather than by individual question:
six conversations form the LoCoMo training pool, and four conversations are
reserved for evaluation. LongMemEval is not used for GRPO training; all 469
filtered LongMemEval questions are used for zero-shot evaluation. For
\bench{}, 99 filtered records are sampled for training, and the remaining 275
records are used for evaluation. The evaluation partitions are not used for
checkpoint selection.

\begin{table}[H]
  \centering
  \caption{Training and evaluation split used for GRPO training and reporting.}
  \label{tab:training-evaluation-split}
  \scriptsize
  \resizebox{\textwidth}{!}{%
  \begin{tabular}{p{0.24\textwidth}rrrrp{0.24\textwidth}}
    \hline\hline
    Data stage & LoCoMo & LongMemEval & \bench{} & Total & Use \\
    \hline
    Original questions & 1,986 & 500 & 500 & 2,986 & Source pool \\
    Filtered questions & 1,536 & 469 & 374 & 2,379 & Candidate pool \\
    Training & 401 & 0 & 99 & 500 & GRPO training \\
    Evaluation & 655 & 469 & 275 & 1,399 & Reporting only \\
    \hline\hline
  \end{tabular}
  }
\end{table}

Each training example stores lightweight supervision and metadata: the user
query, gold answer, gold evidence sessions, source identifier, user or speaker
metadata, source bucket, difficulty label, and evaluation mode. It does not
store fixed retrieved memories or precomputed rollout traces. During GRPO, the
sample metadata is used to call the corresponding online memory-search
environment, retrieve sessions, synthesize events, produce verdicts, and score
the final answer. The policy is therefore trained on online memory-search
behavior rather than on a static retrieved context.

\section{Model and Baseline Configurations}
\label{app:model-combo}

Table~\ref{tab:model-combo} reports the model configuration used by each method
in the main result tables. Every method has a model that produces the final
answer. The distinction is whether the memory component returns evidence to a
separate downstream answerer, as \sys{}, Mem0, and A-Mem do, or whether the
same model directly produces the answer as part of the method, as in
full-context prompting and direct-answerer retrieval agents. GPT-5.1 is used as
the fixed judge for all methods and is not the answerer unless the row
explicitly names GPT-5.1 as the final-answer model.

\begin{table}[H]
  \centering
  \caption{Model combinations used in the main evaluation tables. The
  final-answer model is the model whose output is judged for answer accuracy.
  GPT-5.1 is the fixed judge for all methods, not the answerer for all methods.}
  \label{tab:model-combo}
  \scriptsize
  \resizebox{\textwidth}{!}{%
  \begin{tabular}{p{0.14\textwidth}p{0.17\textwidth}p{0.18\textwidth}p{0.25\textwidth}p{0.18\textwidth}}
    \hline\hline
    Group & Method & Output mode & Memory/search or policy model & Final-answer model \\
    \hline
    Full context & GPT-5.1 & Direct answer & Full history prompt & GPT-5.1 \\
    Full context & Qwen3.5-9B & Direct answer & Full history prompt & Qwen3.5-9B \\
    Write-time memory & Mem0 & Evidence provider & Qwen3.5-9B & Qwen3.5-9B \\
    Write-time memory & A-Mem & Evidence provider & Qwen3.5-9B & Qwen3.5-9B \\
    Search-time retrieval & MemR$^3$ & Direct answer & Qwen3.5-9B & Qwen3.5-9B \\
    Search-time retrieval & Search-R1 & Direct answer & Public checkpoint & Public checkpoint \\
    Temporal memory & Memory-T1-Qwen & Direct answer & Qwen3.5-9B & Qwen3.5-9B \\
    Ours & \sys{}-Vanilla & Evidence provider & Qwen3.5-9B & Qwen3.5-9B \\
    Ours & \sys{}-GRPO & Evidence provider & Qwen3.5-9B GRPO & Qwen3.5-9B \\
    \hline\hline
  \end{tabular}
  }
\end{table}

\section{Training Implementation Details}
\label{app:training-hyperparameters}

\textbf{Policy and rollout.}
The trainable policy is the Qwen3.5-9B Synthesizer, trained with
full-parameter GRPO using VERL/HybridFlow \cite{sheng2024hybridflow}. The
Planner, \sys{} retriever, Qwen3.5-9B answerer, GPT-5.1 judge, and reference
policy are frozen environment components. For each prompt, the actor samples
\(K=8\) trajectories with at most three retrieval rounds. At inference time,
\sys{} uses the same maximum of three retrieval rounds and retrieves the top-7
memories in each round. The rollout context allows 49,152 prompt tokens, 2,048
response tokens, and a 65,536-token SGLang model context
\cite{zheng2023sglang}.

\textbf{Optimization.}
We optimize with AdamW using learning rate \(2{\times}10^{-6}\), weight decay
0.01, betas \((0.9,0.999)\), and Adam epsilon \(10^{-8}\). The learning-rate
schedule is constant with a 3\% warmup, corresponding to 9 warmup steps over
300 global steps. Each GRPO update uses one prompt and eight rollout
trajectories, with PPO mini-batch size 1 and micro-batch size 1. We use PPO
clipping ratio 0.2, actor KL loss coefficient 0.001 with low-variance KL,
gradient clipping 1.0, bf16 training, and FSDP parameter and optimizer
offload. The final run resumes from global step 170 and continues to step 300,
saving checkpoints every 10 steps. Training is conducted on four NVIDIA GH200
96GB GPUs.

\textbf{Reward setting.}
The final run uses reward weights
\[
  (w_o,w_v,w_l)=(0.50,0.10,0.05),
  \qquad c=-0.20.
\]
Appendix~\ref{app:reward-design} defines the reward components, and
Appendix~\ref{app:reward-design-coefficients} analyzes the coefficient choice.

\section{Evaluation Protocol}
\label{app:evaluation-protocol}

The main metric is answer accuracy. In all experiments, GPT-5.1 is used as the
fixed LLM judge for automatic evaluation, and the judge prompt is kept
unchanged across compared methods. BLEU and token-level F1 are not used as
primary metrics because many long-term memory answers are semantically
equivalent despite surface variation, especially for dates, event descriptions,
and short natural-language answers.

LoCoMo is a two-speaker conversational memory benchmark. The evaluation follows
the conversation-memory setting and searches the relevant speaker histories as
required by each question. LongMemEval and \bench{} are treated as single-user
query settings; each query is evaluated independently over its associated
history or generated session set.

Token counts are output-side compactness metrics and are not identical to full
serving cost. For direct-answerer methods, tokens denote final answer-output
tokens. For evidence-provider methods, tokens denote the retrieved evidence or
memory context returned downstream. \sys{}-specific diagnostics additionally
include retrieval rounds and verdict accuracy. Retrieval rounds measure how
often the loop is used, while verdict accuracy measures whether the
Synthesizer's \texttt{sufficient}/\texttt{insufficient} decision agrees with
gold evidence coverage.

\section{Baseline Implementation Notes}
\label{app:baseline-implementation-notes}

This section records implementation adaptations and caveats not captured by
Table~\ref{tab:model-combo}. For stateful memory systems, users and benchmark
files are isolated to avoid cross-example leakage. Local Qwen baselines use
SGLang \cite{zheng2023sglang} with thinking disabled. Embedding-based
baselines use \texttt{text-embedding-3-large} with hidden dimension 3072 unless
otherwise noted. The intended evidence retrieval budget is top-\(k=7\), except
for method-specific budgets or implementation caveats below.

\textbf{Full context.}
Full context does not call a memory system. LoCoMo uses each conversation file
as one full-history prefix. LongMemEval uses each query-specific haystack, and
\bench{} provides all sessions in a sample as context. This is a strong
reference baseline but not a deployable memory system.

\textbf{Mem0.}
We run a local OSS Mem0 codebase \cite{chhikara2025mem0} with lightweight
server/client adaptations. Histories are ingested at the session level, and
Mem0 internally extracts fact-style memory notes; search therefore returns
Mem0 notes rather than raw transcripts. The local backend uses pgvector without
graph-store retrieval and adds synchronous ingestion, JSON repair and retry,
custom fact/update prompts, and top-\(k\) search fixes required by the local
evaluation harness.

\textbf{A-Mem.}
We evaluate A-Mem \cite{xu2025amem} with the released codebase. Each session
is inserted through one \texttt{add\_note()} call with a timestamped header, so
the stored \texttt{MemoryNote.content} is the full transcript rather than an
extracted atomic fact. The vector store is ChromaDB, and the wrapper keeps the
released add/search/evolution path intact. A-Mem evolution updates tags,
context, and links, but does not rewrite the stored transcript into compact
facts; its evidence token counts therefore reflect transcript-level evidence.

\textbf{MemR\(^3\).}
We evaluate MemR\(^3\) \cite{du2025memr3} with the released codebase and its
original retrieve/generate/reflect/answer workflow. Each benchmark item is
converted into MemR\(^3\)'s conversation-plus-QA format: the conversation is
flattened as timestamped speaker text, chunked, embedded, and queried by the
multi-step retrieval loop. The configuration uses chunk size 256, top-\(k=7\),
maximum 5 iterations, and no reranker. We do not train MemR\(^3\). Disabling
the reranker may understate its best retrieval performance, and its score
reflects retrieval plus answer generation rather than pure retrieval accuracy.

\textbf{Search-R1.}
We evaluate Search-R1 \cite{jin2025searchr1} with the public Hugging Face
checkpoint \texttt{PeterJinGo/SearchR1-nq\_hotpotqa\_train-qwen2.5-7b-em-ppo}.
The model is used as a direct answerer with the original
\texttt{<think>}/\texttt{<search>}/\texttt{<answer>} protocol. It can issue up
to three search calls, receives retrieved memory text in \texttt{<information>}
blocks, and the final text inside \texttt{<answer>} is used as the benchmark
response. The checkpoint was trained for open-domain search QA rather than
timestamped long-term conversational memory; the adapter does not fine-tune it
or extract atomic memory facts during ingestion.

\textbf{Memory-T1-Qwen.}
We do not run an official trained Memory-T1 checkpoint because public weights
were not available for the trained model at the time of the experiments. We
therefore report a Memory-T1-style baseline using Qwen3.5-9B rather than
official Memory-T1 checkpoint performance. Session-level SimpleBM25 retrieves
candidate session transcripts; the prompt asks Qwen to select memory IDs and
produce a JSON answer. The retrieval budget is top-\(k=7\), the prompt budget
is 15,500 tokens, temperature is 0, and malformed JSON outputs receive at most
two retries. Session-level BM25 can miss one endpoint of a temporal bridge, and
Qwen3.5-9B is not trained for Memory-T1-style long-term temporal dialogue
reasoning.

\chapter{Algorithm and Reward Details}
\label{app:algorithm-reward-details}

\section{Algorithm Details}
\subsection{NanoMem Training Procedure}
\label{app:nanomem-training-procedure}

This subsection summarizes the online GRPO training loop used to optimize only
the Synthesizer/Compressor policy in \sys. Planner, preprocessor, retriever,
answerer, judge, and reference policy are frozen throughout training.

\begin{tcolorbox}[
  enhanced jigsaw,
  breakable,
  colback=white,
  colframe=black,
  boxrule=0pt,
  borderline north={0.7pt}{0pt}{black},
  borderline south={0.7pt}{0pt}{black},
  sharp corners,
  left=0pt,
  right=0pt,
  top=2pt,
  bottom=2pt,
  before skip=0.6em,
  after skip=0.8em,
]
\scriptsize
\noindent\textbf{Algorithm 1} GRPO Training for Multi-turn Memory Evidence Synthesis

\vspace{0.15em}
\hrule
\vspace{0.45em}

\noindent\textbf{Input:}\\
\hspace*{1.5em}Training set $\mathcal{D}_{\mathrm{train}}$; each sample contains query $q$, gold answer $a^\star$,\\
\hspace*{1.5em}gold evidence sessions $\mathcal{E}^\star$, source DB path, user id, and analysis metadata\\
\hspace*{1.5em}Trainable compressor policy $\pi_\theta$ initialized from Qwen3.5-9B\\
\hspace*{1.5em}Frozen planner, retriever, answerer, judge, and reference policy $\pi_{\mathrm{ref}}$\\
\hspace*{1.5em}Rollout group size $G=8$, maximum search rounds $R_{\max}=3$

\vspace{0.25em}
\noindent\textbf{Ensure:}\\
\hspace*{1.5em}Optimized compressor policy $\pi_{\theta'}$

\vspace{0.35em}
\begin{tabular}{@{}r@{\quad}p{0.84\linewidth}@{}}
1: & \textbf{for each} training step \textbf{do} \\
2: & \hspace*{1em}Sample one training example $x=(q,a^\star,\mathcal{E}^\star,\mathrm{DB},u,\mathrm{meta})$ from $\mathcal{D}_{\mathrm{train}}$ \\
3: & \hspace*{1em}\textbf{Stage 1: Input and initialization} \\
4: & \hspace*{2em}Load the sample-specific memory DB and user id; keep planner, retriever, answerer, judge, and $\pi_{\mathrm{ref}}$ frozen \\
5: & \hspace*{1em}\textbf{Stage 2--3: Online multi-turn group rollout} \\
6: & \hspace*{2em}\textbf{for} rollout $g=1,\ldots,G$ \textbf{do} \\
7: & \hspace*{3em}Initialize previous events $\mathcal{M}^{(0)}\leftarrow\emptyset$, previous cues $\mathcal{C}^{(0)}\leftarrow\emptyset$, and terminal flag $\mathrm{stop}\leftarrow\mathrm{false}$ \\
8: & \hspace*{3em}\textbf{for} round $r=1,\ldots,R_{\max}$ \textbf{do} \\
9: & \hspace*{4em}Use the frozen planner to generate retrieval cues $\mathcal{C}^{(r)}$ from $q$, $\mathcal{C}^{(r-1)}$, and compressor reasoning when available \\
10: & \hspace*{4em}Use the frozen retriever to search sessions $\mathcal{S}^{(r)}$ from the sample-specific memory DB \\
11: & \hspace*{4em}Construct temporal hints online from the query and retrieved sessions \\
12: & \hspace*{4em}Sample compressor XML $z_g^{(r)}\sim\pi_\theta(\cdot\mid q,\mathrm{hints},\mathcal{M}^{(r-1)},\mathcal{S}^{(r)})$ \\
13: & \hspace*{4em}Parse $z_g^{(r)}$ into reasoning, events $\mathcal{M}^{(r)}$, and verdict $v^{(r)}$ \\
14: & \hspace*{4em}\textbf{if} $v^{(r)}=\mathrm{sufficient}$, no new evidence is found, or XML is invalid \textbf{then} set $\mathrm{stop}\leftarrow\mathrm{true}$ and break \\
15: & \hspace*{3em}\textbf{end for} \\
16: & \hspace*{3em}Use the frozen answerer to generate final answer $\hat{a}_g$ from $q$ and terminal events $\mathcal{M}^{(r)}$ \\
17: & \hspace*{3em}Store rollout trace $\tau_g$: retrieved sessions, compressor XML, reasoning, events, verdicts, terminal reason, and answer \\
18: & \hspace*{2em}\textbf{end for} \\
19: & \hspace*{1em}\textbf{Stage 4: Reward computation} \\
20: & \hspace*{2em}\textbf{for} rollout $g=1,\ldots,G$ \textbf{do} \\
21: & \hspace*{3em}Compute format reward from XML validity and schema compliance \\
22: & \hspace*{3em}Compute verdict reward by comparing sufficient/insufficient decisions with gold evidence coverage $\mathcal{E}^\star$ \\
23: & \hspace*{3em}Use the frozen judge to compare $\hat{a}_g$ with $a^\star$ and assign outcome reward \\
24: & \hspace*{3em}Compute length reward from compressor output compactness, mainly active when the outcome is correct \\
25: & \hspace*{3em}Combine format, verdict, outcome, and length rewards into scalar reward $R_g$ \\
26: & \hspace*{2em}\textbf{end for} \\
27: & \hspace*{1em}\textbf{Stage 5: GRPO policy update} \\
28: & \hspace*{2em}Compute group-relative advantages from $\{R_g\}_{g=1}^{G}$ for the same query \\
29: & \hspace*{2em}Increase likelihood of higher-reward compressor outputs and decrease likelihood of lower-reward outputs \\
30: & \hspace*{2em}Apply KL regularization against frozen reference policy $\pi_{\mathrm{ref}}$ \\
31: & \hspace*{2em}Update only compressor policy parameters $\theta$ \\
32: & \textbf{end for} \\
33: & Select the final checkpoint using training reward only; held-out sets are used for reporting \\
34: & \textbf{return} optimized compressor policy $\pi_{\theta'}$
\end{tabular}
\end{tcolorbox}

\subsection{Reward Design}
\label{app:reward-design}

We define the reward over a complete rollout trajectory
$\tau=\{(h_1,z_1),\ldots,(h_T,z_T)\}$, where $h_t$ is the retrieval context and
$z_t=(r_t,\mathcal{E}_t,\hat{v}_t)$ is the Synthesizer action at round $t$. For
coefficient analysis, the same trajectory is abstracted as
$\tau=(f,T,\mathbf{v},o,\ell)$, where $f$ is format validity, $T$ is the number
of retrieval rounds, $\mathbf{v}$ is the sequence of verdict correctness
signals, $o$ is final answer correctness, and $\ell$ is evidence compactness.

\noindent\textbf{Format reward.}
Let $f(\tau)\in\{0,1\}$ indicate whether all Synthesizer outputs in the rollout
are parseable under the required XML schema. Format validity acts as a gate:
if $f(\tau)=0$, all other reward terms are not evaluated and the trajectory
receives the fixed reward $c<0$. Equivalently, the format-failure penalty
magnitude is $-c>0$.

\noindent\textbf{Length reward.}
Let $m(z_t)$ be the token length of the Synthesizer output at round $t$.
The compactness score $\ell(\tau)\in[0,1]$ is activated only when the final
answer is correct:
\[
\ell(\tau)
=
\mathbb{I}\{R_{\mathrm{out}}(\tau)=1\}
\cdot
\mathrm{clip}\!\left(
1-\frac{1}{Tm_{\max}}\sum_{t=1}^{T}m(z_t),
0,1
\right),
\]
where $m_{\max}$ is the length threshold. This term prefers shorter evidence
among trajectories that already solve the task.

\noindent\textbf{Verdict reward.}
The Synthesizer must decide whether the current retrieved evidence covers the
query or whether a missing evidence gap remains. Let $\mathcal{G}_q$ denote the
gold evidence session set for query $q$, and let
$\mathcal{S}^{\mathrm{ret}}_{\le t}$ denote the session identifiers covered by
the accumulated retrieved sessions $\widetilde{\mathcal{S}}_{\le t}$. The
round-level verdict reward is
\[
v_t=
\begin{cases}
+1, & \mathcal{G}_q\subseteq\mathcal{S}^{\mathrm{ret}}_{\le t}
\ \text{and}\ \hat{v}_t=\texttt{sufficient},\\
+1, & \mathcal{G}_q\not\subseteq\mathcal{S}^{\mathrm{ret}}_{\le t}
\ \text{and}\ \hat{v}_t=\texttt{insufficient},\\
-1, & \text{otherwise}.
\end{cases}
\]
This term evaluates whether the policy correctly decides to stop or continue
under the objective retrieval state. Since it is computed independently at each
round, the trajectory-level verdict score is $V=\sum_{t=1}^{T}v_t$, providing
fine-grained process supervision for multi-round search.

\noindent\textbf{Outcome reward.}
The outcome reward ties the final Temporal Evidence Pool to task success. Let
$H_{T+1}$ be the Temporal Evidence Pool after the final round. The frozen answer
model generates $\hat{a}=\mathrm{Ans}(q,H_{T+1})$, and the evaluator checks
semantic equivalence with the gold answer $a^\star$:
\[
o(\tau)=\mathbb{I}\{\mathrm{Eval}(\hat{a},a^\star)=1\}.
\]
This component measures whether the constructed evidence state is sufficient
for answering the original query.

Combining these components, the total reward is
\[
R(\tau)
=
\begin{cases}
c, & f(\tau)=0,\\
w_o o(\tau)+w_v\sum_{t=1}^{T}v_t+w_l\ell(\tau), & f(\tau)=1,
\end{cases}
\]
where $w_o,w_v,w_l>0$ are the weights for outcome, verdict, and compactness.
Thus, $o$ and $\ell$ are trajectory-level signals, while $v_t$ is accumulated
over retrieval rounds and format validity gates the whole trajectory. We do not
introduce a separate event-time reward: incorrect event-time reasoning is
reflected in $o(\tau)$ through retrieval drift and final answer failure.

\subsection{Reward Coefficient Analysis}
\label{app:reward-design-coefficients}

\noindent\textbf{Problem Definition.}
We abstract a rollout trajectory as
$\tau=(f,T,\mathbf{v},o,\ell)$, where $f\in\{0,1\}$ denotes whether the
Synthesizer output is parseable under the required XML schema,
$T\in\{1,\ldots,T_{\max}\}$ is the number of retrieval rounds, and
$\mathbf{v}=(v_1,\ldots,v_T)$ records round-level verdict correctness with
$v_t\in\{-1,+1\}$. The final answer correctness is $o\in\{0,1\}$, and
$\ell\in[0,1]$ measures evidence compactness, with $\ell=0$ when $o=0$ because
compactness is only rewarded for correct outcomes. Let
$V=\sum_{t=1}^{T}v_t$ be the cumulative verdict score. Format acts as a gate:
if $f=0$, verdict, outcome, and length terms are not evaluated and the rollout
receives a fixed penalty $c$. Otherwise, the reward is:
\[
R(\tau)=
\begin{cases}
c, & f=0,\\
w_vV+w_o o+w_l\ell, & f=1,
\end{cases}
\]
where $c\in\mathbb{R}$ is the format-failure penalty and
$w_v,w_o,w_l>0$ are reward weights for verdict correctness, answer correctness,
and compactness, respectively. With $T_{\max}=3$, the wrong-outcome reward range
under format success is $[-3w_v,3w_v]$, while the correct-outcome reward range
is $[-3w_v+w_o,3w_v+w_o+w_l]$.

We require this reward to satisfy six ordering goals. (G1) \textit{Format
gate}: every format-valid trajectory should exceed format failure. (G2)
\textit{Outcome dominance}: any correct-outcome trajectory should exceed any
wrong-outcome trajectory, regardless of verdict history or compactness. (G3)
\textit{Verdict discrimination}: for fixed outcome and compactness, a larger
cumulative verdict score should receive larger reward. (G4) \textit{Multi-round
credit}: a trajectory with more correct verdict rounds should receive more
verdict credit. (G5) \textit{Compactness preference}: for fixed outcome and
verdict, greater compactness should receive larger reward. (G6)
\textit{Verdict over length}: a single verdict correction should dominate the
full compactness range, so process supervision cannot be overridden by the
length term.

We also define four guaranteed margins to express the intended priority among
learning signals. $\Delta_1$ is the \emph{format-gate margin}: it asks how far
the worst format-success trajectory remains above format failure. Format
failure receives the fixed reward $c$, while the worst parseable trajectory has
all other signals minimized, namely $V=-3$, $o=0$, and $\ell=0$:
\[
R^{\inf}_{f=1}=w_v(-3)+w_o\cdot 0+w_l\cdot 0=-3w_v,
\]
\[
\Delta_1=R^{\inf}_{f=1}-c=-3w_v-c.
\]
Thus, $\Delta_1>0$ guarantees that any parseable output is preferred to a
format failure. $\Delta_2$ is the \emph{outcome-discrimination margin}: under
format success, it compares the worst correct-answer trajectory with the best
wrong-answer trajectory. The former has $V=-3$ and $\ell=0$, while the latter
has $V=3$:
\[
R^{\inf}_{o=1}=-3w_v+w_o,\quad
R^{\sup}_{o=0}=3w_v,
\]
\[
\Delta_2=R^{\inf}_{o=1}-R^{\sup}_{o=0}=w_o-6w_v.
\]
The term $6w_v$ is the maximal swing introduced by cumulative verdict scoring;
$\Delta_2>0$ means that answering correctly remains better than answering
incorrectly even under this extreme verdict contrast. $\Delta_3$ is the
\emph{single-step verdict signal}: changing one round-level verdict from
incorrect to correct changes $v_t$ from $-1$ to $+1$, increasing $V$ by $2$:
\[
w_v(V+2)-w_vV=2w_v,
\]
\[
\Delta_3=2w_v.
\]
This margin is independent of $T$, $o$, and $\ell$. Finally, $\Delta_4$ is the
\emph{compactness range}: since $\ell\in[0,1]$ and the length term is
$w_l\ell$, the maximum reward variation induced by compactness is
\[
w_l\cdot 1-w_l\cdot 0=w_l,
\]
\[
\Delta_4=w_l.
\]
We require $\Delta_1>\Delta_2>\Delta_3>\Delta_4>0$ so that the guaranteed
discrimination decreases along the intended priority order: format validity,
answer correctness, verdict quality, and evidence compactness. This ordering
ensures that GRPO receives policy-gradient signals aligned with the desired
learning hierarchy.

\noindent\textbf{Constraint Derivation.}
We first derive constraints directly from the six design goals.

\noindent\emph{C1
(Format gate).} The worst format-valid trajectory has $o=0$ and
$V=-T_{\max}=-3$, so Goal 1 requires it to beat the format-failure penalty:
\[
-3w_v>c.
\]
\emph{C2 (Outcome dominance).} The worst correct trajectory must beat the
best wrong trajectory:
\[
R^{\inf}_{o=1}>R^{\sup}_{o=0}
\quad\Longrightarrow\quad
-3w_v+w_o>3w_v
\quad\Longrightarrow\quad
w_o>6w_v.
\]
\emph{C3 (Verdict discrimination and multi-round credit).} Since the reward
is linear in $V$ with slope $w_v$, verdict ordering and cumulative multi-round
credit require
\[
w_v>0.
\]
\emph{C4 (Compactness preference).} Since compactness only affects correct
trajectories through $w_l\ell$, compactness preference requires
\[
w_l>0.
\]
\emph{C5 (Verdict over length).} A single verdict correction changes
$V$ by $2$, so it must dominate the entire compactness range:
\[
2w_v>w_l.
\]

\noindent\textbf{Gap-Ordering Constraints.}
We next derive constraints from the desired hierarchy
$\Delta_1>\Delta_2>\Delta_3>\Delta_4>0$.

\noindent\emph{G1} follows from
$\Delta_1>\Delta_2$:
\[
-3w_v-c>w_o-6w_v
\quad\Longrightarrow\quad
c<3w_v-w_o.
\]
\emph{G2} follows from $\Delta_2>\Delta_3$:
\[
w_o-6w_v>2w_v
\quad\Longrightarrow\quad
w_o>8w_v.
\]
\emph{G3} follows from $\Delta_3>\Delta_4$:
\[
2w_v>w_l.
\]
\emph{G4} follows from $\Delta_4>0$:
\[
w_l>0.
\]

Combining the six ordering goals with the gap hierarchy gives the feasible
coefficient region:
\[
\mathcal{W}=
\left\{(c,w_v,w_o,w_l)\in\mathbb{R}\times\mathbb{R}_{>0}^{3}\;\middle|\;
\begin{aligned}
&w_v>0,\\
&w_o>8w_v,\\
&c<3w_v-w_o,\\
&2w_v>w_l
\end{aligned}
\right\}.
\]
This region keeps outcome correctness as the dominant task-level signal,
preserves verdict correctness as process-level supervision, and treats format
and length terms as stabilizing constraints rather than primary task objectives.

\noindent\textbf{Coefficient Selection.}
We select coefficients in three steps. First, we instantiate the design goals
as a four-level priority structure: answer correctness is the primary objective,
verdict correctness provides process-level control over stopping and continued
search, compactness improves efficiency, and format validity acts as a hard
gate. Second, this priority induces ratio constraints among weights: $w_o$
should dominate $w_v$, while the single-step verdict signal $2w_v$ should cover
the full length range $w_l$, preventing the policy from sacrificing verdict
quality for shorter evidence. Third, we choose the default setting
$(w_o,w_v,w_l)=(0.50,0.10,0.05)$. This makes outcome the dominant signal in our
rollout distribution, gives a verdict-step margin $2w_v=0.20$ larger than the
length range $w_l=0.05$, and sets the format penalty by
$c=3w_v-w_o=-0.20$ so that format failure remains below the worst valid output.

\noindent\textbf{Effectiveness Analysis.}
Our construction can be viewed as an expressibility argument for a desired
trajectory ordering. Following the perspective of
Abel et al.~\citep{abel2021expressivity}, reward design is not merely assigning
heuristic scalar scores: it asks whether a reward function can express a
task-level preference relation. In our case, the six design goals define a
partial order over rollout trajectories for temporal-causal evidence synthesis.
For example, outcome dominance requires every correct-answer trajectory to
outrank every wrong-answer trajectory, while the verdict and length goals impose
finer preferences within trajectory subsets. The feasible region $\mathcal{W}$
then gives a constructive certificate that this ordering is expressible by our
multi-level reward family. Moreover, the gap hierarchy
$\Delta_1>\Delta_2>\Delta_3>\Delta_4>0$ strengthens strict preference by
requiring larger guaranteed margins for higher-priority distinctions, echoing
the range-based view that acceptable behaviors should be separated from
unacceptable ones by reward margins. This explains why the coefficient design is
effective: it aligns the scalar GRPO learning signal with the intended
trajectory-level preference hierarchy, rather than relying on unstructured
reward mixing.


\chapter{Dataset Details}
\label{app:dataset-details}

\section{Dataset Construction}
\subsection{TimeMemEval Benchmark Construction}
  \label{app:timememeval-construction}

  \begin{tcolorbox}[
    enhanced jigsaw,
    breakable,
    colback=white,
    colframe=black,
    boxrule=0pt,
    borderline north={0.7pt}{0pt}{black},
    borderline south={0.7pt}{0pt}{black},
    sharp corners,
    left=0pt,
    right=0pt,
    top=2pt,
    bottom=2pt,
    before skip=0.6em,
    after skip=0.8em,
  ]
  \scriptsize
  \noindent\textbf{Algorithm 2} \textsc{TimeMemEval} Benchmark Construction Procedure

  \vspace{0.15em}
  \hrule
  \vspace{0.45em}

  \noindent\textbf{Require:}\\
  \hspace*{1.5em}Persona/theme pools, event-family library, bridge scopes
  $\mathcal{B}$, offset sampler $\Omega$, terminal types $\mathcal{Y}$\\
  \hspace*{1.5em}Distractor modes $\mathcal{M}$, surface temporal styles $\mathcal{R}$,
  event/session generators, uniqueness judge $\mathcal{J}$, temporal validator
  $\mathcal{V}$, retry budgets

  \vspace{0.25em}
  \noindent\textbf{Ensure:}\\
  \hspace*{1.5em}\textsc{TimeMemEval} records with \texttt{record.json},
  \texttt{sessions.json}, generation traces, split metadata, and validation logs

  \vspace{0.25em}
  \begin{tabular}{@{}r@{\quad}p{0.84\linewidth}@{}}
  1: & \textbf{function} BuildTimeMemEval$(\mathcal{B},\Omega,\mathcal{Y})$ \\
  2: & \hspace*{1em}\textbf{for each} requested record \textbf{do} \\
  3: & \hspace*{2em}Sample persona/theme context and generate a source event $e_s$ with a valid temporal anchor \\
  4: & \hspace*{2em}Sample bridge scope and offset $(\beta,\delta)$, where $\beta\in\mathcal{B}$ and $\delta\in\Omega$ \\
  5: & \hspace*{2em}$W^\star \leftarrow$ ComputeTargetWindow$(e_s.\mathrm{event\_time},\beta,\delta)$ \\
  6: & \hspace*{2em}Sample terminal type $y\in\mathcal{Y}$ and a compatible event family; generate destination event $e_d$ inside $W^\star$ \\
  7: & \hspace*{2em}Generate probe question $q$ and answer $a$ that require linking $e_s$ to $e_d$ through $W^\star$ \\
  8: & \hspace*{2em}Generate hard distractors $\mathcal{D}$ using modes in $\mathcal{M}$, including same-activity, answer-alias, near-window, and source-anchor-alias distractors outside $W^\star$ \\
  9: & \hspace*{2em}Generate filler events $\mathcal{U}$ that add dialogue noise without leaking the answer or temporal bridge \\
  10: & \hspace*{2em}\textbf{if} $\mathcal{J}(e_s,e_d,\mathcal{D},\mathcal{U},q,a)$ rejects uniqueness or answerability \textbf{then} retry the failed stage \\
  11: & \hspace*{2em}\textbf{for each} event $e\in\{e_s,e_d\}\cup\mathcal{D}\cup\mathcal{U}$ \textbf{do} \\
  12: & \hspace*{3em}Sample a surface temporal style $r\in\mathcal{R}$ and generate a multi-turn session $s_e$ with relative, indirect, or calendar-based temporal phrasing \\
  13: & \hspace*{3em}\textbf{if} $\mathcal{V}(s_e,e,\beta)$ rejects temporal consistency \textbf{then} regenerate $s_e$ or its dependent event \\
  14: & \hspace*{2em}\textbf{end for} \\
  15: & \hspace*{2em}Write record, sessions, traces, split metadata, and validation logs \\
  16: & \hspace*{1em}\textbf{end for} \\
  17: & \textbf{end function}
  \end{tabular}
  \end{tcolorbox}

  \textsc{TimeMemEval} is designed as a controlled benchmark for temporal-bridge
  retrieval in long-term memory, rather than covering all forms of
  temporal reasoning. Each instance requires a model to identify a source event,
  derive a target temporal window, retrieve a destination event in that window,
  and answer a terminal question about the destination event. This construction
  intentionally isolates interval-conditioned multi-hop retrieval from ordinary
  single-hop lookup, while keeping explicit metadata that allows us to analyze the
  effect of bridge scope, offset, terminal type, distractor type, and surface
  temporal phrasing.

  To reduce benchmark artifacts, generation is not tied to a single template.
  Records are sampled from persona/theme contexts, event-family libraries,
  terminal-question types, bridge scopes, offset ranges, and multiple distractor
  modes. Distractors are constructed to be semantically and structurally close to
  the destination event: for example, they may share the same activity, preserve
  the same object or location skeleton when required by the terminal type, occur
  near the target window, or introduce competing source-like temporal anchors.
  The final judge enforces that the question is not answerable from the source
  alone or the destination alone, but remains uniquely answerable when all source,
  destination, distractor, and filler sessions are present.

  Session generation is separated from event generation. After the event-level
  graph passes the uniqueness gate, each event is realized as a multi-turn memory
  session with a sampled temporal surface form. These forms include deterministic
  calendar references, relative references such as ``last week'' or ``three months
  ago'', and less calendar-like user-memory anchors such as life episodes,
  projects, trips, moves, medical events, habits, or emotionally salient periods.
  This is intended to make the benchmark less dependent on explicit date strings:
  the temporal validator checks whether the realized session still supports the
  intended day, week, or month bucket under the sampled bridge scope.

  We report \textsc{TimeMemEval} as a targeted temporal-bridge retrieval stress
  test. It does not subsume broader temporal phenomena such as long causal chains,
  belief revision, recurring habits, partial temporal order, or contradictory
  memories. To make this limitation explicit, benchmark metadata records the
  temporal topology of every instance, and evaluation can be stratified by bridge
  scope, terminal type, distractor mode, surface temporal style, and whether
  temporal normalization is provided to the model. This allows us to distinguish
  performance gains from temporal normalization, retrieval policy, and reasoning
  over the generated memory evidence, and supports held-out splits over event
  families, surface templates, and temporal relation types.

\subsection{Training and Held-out Evaluation Dataset Construction}
\label{app:train-valid-data-construction}

We construct the GRPO data from two public long-term memory benchmarks,
LoCoMo and LongMemEval, together with our generated \bench records. LoCoMo and
LongMemEval provide real multi-session dialogue questions with gold answers and
evidence-session annotations, while \bench contributes targeted
temporal-causal questions whose answer evidence is not directly addressable from
the original query. These sources play different roles: the public benchmarks
preserve general long-memory coverage, whereas \bench increases the density of
training and evaluation cases that require temporal bridge construction.

\textbf{Source filtering.}
We start from 1,986 LoCoMo questions, 500 LongMemEval questions, and 500 raw
\bench candidates. For LoCoMo, we exclude category-5 adversarial questions; for
LongMemEval, we exclude single-session preference questions, yielding 1,536 and
469 usable public questions, respectively. For \bench, the generation pipeline
first produces 500 raw candidates and then applies automatic validity checks for
invalid generation attempts, duplicate examples, split leakage, incomplete
metadata, temporal inconsistency among target windows, session times and event
times, and answer non-uniqueness. This leaves 374 filtered \bench records that
are eligible for training or held-out evaluation. The resulting filtered pool
contains 2,379 examples.

\textbf{Training and held-out evaluation splits.}
The final GRPO training set contains 500 examples. To avoid conversation-level
leakage in LoCoMo, we first partition LoCoMo by dialogue rather than by
question: six conversations form the LoCoMo training pool and the remaining four
conversations are reserved for held-out evaluation. We then sample 401 LoCoMo
questions from the six training conversations. LongMemEval is not used for GRPO
training; all 469 filtered LongMemEval questions are kept for zero-shot
held-out evaluation. For \bench, we sample 99 training examples from the 374
filtered records and use the remaining 275 filtered records for held-out
evaluation. The held-out evaluation sets are not used for checkpoint selection.

\textbf{Example format.}
Each training example stores only lightweight supervision and metadata: the user
query, gold answer, gold evidence sessions, source identifier, user or speaker
metadata, source bucket, difficulty label, and evaluation mode. It does not store
fixed retrieved memories or precomputed rollout traces. During GRPO, the system
uses the sample metadata to dynamically call the corresponding NanoMem search
environment, generate temporal hints, retrieve sessions, synthesize events,
produce verdicts, and score the final answer. This ensures that the policy is
trained on online memory-search behavior rather than a static retrieved context.

\begin{table}[h]
\centering
\footnotesize
\setlength{\tabcolsep}{3pt}
\begin{tabular}{@{}lrrrrp{0.23\linewidth}@{}}
\toprule
\textbf{Data stage} & \textbf{LoCoMo} & \textbf{LongMemEval} &
\textbf{TME} & \textbf{Total} & \textbf{Use} \\
\midrule
Original questions & 1,986 & 500 & 500 & 2,986 & Source pool \\
Filtered questions & 1,536 & 469 & 374 & 2,379 & Candidate pool \\
Train & 401 & 0 & 99 & 500 & GRPO training \\
Held-out evaluation & 655 & 469 & 275 & 1,399 & Reporting only \\
\bottomrule
\end{tabular}
\caption{Distribution of query pools for GRPO training and held-out evaluation.
LoCoMo is split by conversation; LME is zero-shot; TME uses filtered records.}
\label{tab:train-valid-data-distribution}
\end{table}


\chapter{Prompt Templates and Case Study}
\label{app:prompt-templates-and-case-study}

\section{Prompt Templates}
\tcbset{
  appendixexample/.style={
    enhanced jigsaw,
    breakable,
    width=\linewidth,
    colback=black!4,
    colframe=black!60,
    colbacktitle=black!65,
    coltitle=white,
    fonttitle=\small\rmfamily,
    boxrule=0.8pt,
    arc=2pt,
    left=8pt,
    right=8pt,
    top=6pt,
    bottom=6pt,
    before skip=0.45em,
    after skip=0.55em,
  }
}

\subsection{Prompt Template of Initial Planner}
\label{app:planner-initial-prompt}

\begin{tcolorbox}[appendixexample,title={Prompt Template of Initial Planner: $P_{\mathrm{plan}}^{(0)}$}]
\begin{Verbatim}[fontsize=\scriptsize,formatcom=\linespread{0.82}\selectfont,breaklines=true,breakanywhere=true]
You are a memory retrieval planner. Given a user query, decompose it into
targeted retrieval cues for searching the user's conversation history.

## Output format (XML only)

<cues>
  <cue>
    <text>retrieval query text</text>
  </cue>
</cues>

---

## Query

{query}

---

## Rules

Cue text:
- Decompose the query into one or more focused retrieval cues.
- Each cue should target a distinct aspect or entity of the query.
- Always generate at least one `<cue>`. If not sure what to generate,
  copy the original user query.
- Preserve named entities, dates, locations, event types, quantities, and
  temporal semantics.
- Preserve the user's intent as closely as possible. Do not broaden a
  specific question into a generic summary query.
- For questions asking "when", keep the cue as a "when" query about the
  same event.
- For dependent temporal questions, emit the original standalone cue plus
  support cues for the anchor event and target fact.
- For count, comparison, ordering, or duration questions, emit support cues
  that retrieve the underlying items, quantities, event times, or start/end
  anchors.
- Keep the original standalone intent as the first cue, then add supporting
  cues.
- Do not annotate temporal phrases. Temporal extraction is handled by
  deterministic preprocessing after cue generation.

Return XML only, no markdown fences, no text outside the tags.
\end{Verbatim}
\end{tcolorbox}
\captionof{figure}{Prompt template used by the initial planner to decompose the user query into targeted retrieval cues.}

\subsection{Prompt Template of Planner Retry}
\label{app:planner-retry-prompt}

\begin{tcolorbox}[appendixexample,title={Prompt Template of Planner Retry: $P_{\mathrm{plan}}^{(\mathrm{retry})}$}]
\begin{Verbatim}[fontsize=\scriptsize,formatcom=\linespread{0.82}\selectfont,breaklines=true,breakanywhere=true]
You are a memory retrieval planner. A previous retrieval round was
attempted but the synthesizer judged the evidence insufficient.
Your task is to generate new targeted cues that address the identified gap.

## Output format (XML only)

<cues>
  <cue>
    <text>retrieval query text</text>
  </cue>
</cues>

---

## Original query

{query}

## Previous cues

{previous_cues}

## Temporal Evidence Pool
{pool}

## Synthesizer reasoning from previous round

{previous_reason}

---

## Rules

Cue text:
- Generate 1-3 cues that target the missing evidence gap named in the
  synthesizer reasoning and its extracted temporal evidences.
- If reasoning found an anchor/bridge entity/place/time, reuse that entity to
  guide the new cue generation.
- Combine the anchor/bridge with the missing target question relation, e.g.
  added beside, requested from, mailed for, set up through, etc.
- Do not re-search only for the anchor/source if it was already found in reason;
  search for the downstream action/object involving the anchor.
- Turn "no evidence of X near/from/through Y" into a cue for "X
  near/from/through Y". It is your job to recover missing relations through new
  cues.
- Do not repeat previous cues unless the reasoning explicitly asks for
  the same topic with a different temporal anchor or framing.
- Preserve the user's intent as closely as possible. Do not broaden a
  specific question into a generic summary query.
- Keep cues short and search-like; avoid meta questions such as
  "is this linked" or "is this the same item".
- For questions asking "when", keep the cue as a "when" query about the
  same event.
- For count, comparison, ordering, or duration questions, emit support cues
  that retrieve the underlying items, quantities, event times, or start/end
  anchors.
- Do not annotate temporal phrases. Temporal extraction is handled by
  deterministic preprocessing after cue generation.

Return XML only, no markdown fences, no text outside the tags.
\end{Verbatim}
\end{tcolorbox}
\captionof{figure}{Prompt template used by the retry planner to generate new cues from the previous Synthesizer reasoning.}

\subsection{Prompt Template of Synthesizer}
\label{app:synthesizer-prompt}

\begin{tcolorbox}[appendixexample,title={Prompt Template of Synthesizer: $P_{\mathrm{synth}}$}]
\begin{Verbatim}[fontsize=\scriptsize,formatcom=\linespread{0.82}\selectfont,breaklines=true,breakanywhere=true]
You are a memory evidence synthesizer. Given a user query, temporal hints,
previously extracted events, and retrieved conversation sessions, extract
structured evidence and judge whether retrieval is sufficient to answer the
query.

## Output format (XML only, this exact order)

<reasoning>
Concise verdict justification (<=200 words): what decisive evidence is
present, what is missing if any, and why the verdict follows.

</reasoning>

<events>
  <event
    session_id="source session ID; exactly one session per event"
    session_time="YYYY-MM-DDTHH:MM:SS.sssZ"
    event_time="YYYY-MM-DD or 'unknown'">
    A single self-contained evidence statement. Free text, one paragraph,
    no nested tags, no markdown.
  </event>
</events>

<verdict>sufficient|insufficient</verdict>

---

## Normalized Query

{query}

## Retrieved sessions

{sessions}

Retrieved session text preserves original temporal phrases and may annotate
them with parenthesized normalized times, such as `last night (during
2023-05-27 18:00:00 to 2023-05-28 05:59:59)`. These parenthesized values are
deterministic parser aids for interpreting time; they are not additional facts
or evidence beyond the session text.

## Temporal Evidence Pool

{previous_events}

Temporal Evidence Pool are evidence already extracted in earlier retrieval rounds.
They may summarize facts that also appear in the current retrieved sessions.
Use them as context for sufficiency and deduplication, but do not treat them
as new evidence or a substitute for facts grounded in retrieved sessions.
---

## Rules

`Verdict`:
`insufficient` - critical evidence clearly absent; loop continues.
`sufficient` - events reliably cover what the query needs; loop ends.

`event_time`: resolve relative expressions using session timestamps and
temporal hints. Each temporal hint contains the original query phrase and an
absolute UTC timestamp range in `iso_span`. If no explicit date is
recoverable, fall back to the session date. Use `unknown` only if even the
session date is unavailable.

Output rules:
- The final `<events>` output is cumulative. Include all useful
  `<previous_events>` first, unchanged when still valid, then append newly
  extracted events from the current retrieved sessions. Do not delete or
  rewrite previous events unless you find an error that can be corrected by
  current sessions.
- One `<event>` per distinct evidence item, grounded in exactly one session
  (`session_id` must be a single session ID from the input; one session may
  yield multiple events). Always emit partial events even when the verdict is
  `insufficient`. Use `<events/>` only if no useful evidence exists at all.
- If a fact is already covered by previous events and the current sessions add
  no new information, do not repeat it as a new event.
- Preserve temporal relations (before/after/during, duration anchors).
- Event text: free text, one paragraph, no nested tags, no markdown.
- Event text should avoid including temporal phrases when possible; represent
  time in the `event_time` attribute instead.
- Return XML only, no text outside the tags.
- All three attributes (`session_id`, `session_time`, `event_time`) are
  required on every `<event>`.
- For quantity, state, or collection queries, retain all update events
  in chronological order; do not drop later increments or corrections.

\end{Verbatim}
\end{tcolorbox}
\captionof{figure}{Prompt template used by the Synthesizer to extract temporal evidence and decide whether retrieval is sufficient.}

\subsection{Case Study: Temporal Bridge Retrieval}
\label{app:case-study-temporal-bridge}

\begin{tcolorbox}[appendixexample,title={Case Study: Temporal Bridge Retrieval}]
\textbf{Case overview.}
This case illustrates how \sys recovers an answer when the decisive temporal
anchor is not stated in the original query but must be inferred from a bridge
session retrieved in an earlier round.

\smallskip
\textbf{Query.}
What did Haruki read 4 months after sorting through old notebooks and setting
aside pages that could become personal essays?

\smallskip
\textbf{Round 1.}
The initial planner issues cues for the original question, the notebook-sorting
event, and Haruki's reading activities. Retrieval returns several reading
distractors from February, April, May, September, November, and December,
together with the bridge session \texttt{S\_E001}. The key bridge evidence is:
\emph{Haruki: I spent most of the afternoon on the living room floor with old
notebooks and binders, and 2 months ago I made a separate pile of pages that
could become personal essays instead of project records.}

\smallskip
Since \texttt{S\_E001} is timestamped 2024-04-18, the normalized phrase ``2
months ago'' grounds the sorting event in February 2024. The Synthesizer then
infers that four months after February 2024 is June 2024. However, because the
retrieved sessions contain no June reading evidence, it returns
\texttt{insufficient}.

\smallskip
\textbf{Retry cues.}
\begin{Verbatim}[fontsize=\scriptsize,formatcom=\linespread{0.82}\selectfont,breaklines=true,breakanywhere=true]
What did Haruki read in June 2024 (during 2024-06)?
What books did Haruki read between May 2024 and September 2024?
What reading activities did Haruki have in June 2024 (during 2024-06)?
\end{Verbatim}

\smallskip
\textbf{Round 2.}
The second retrieval round finds the target session
\texttt{S\_H2\_E001\_01}, timestamped 2024-06-08:
\emph{Haruki: I read The Death and Life of Great American Cities this month,
and it got me thinking about how to turn my old engineering observations into
personal essays instead of reports.}

\smallskip
\textbf{Final synthesis.}
The final Synthesizer output carries forward the bridge event, appends the June
reading event, and marks the evidence as \texttt{sufficient}. The final answer
is \emph{The Death and Life of Great American Cities}, matching the ground-truth
label. This example highlights the role of Synthesizer reasoning as the
information bridge between rounds: it converts first-round bridge evidence into
a concrete temporal cue that the retry planner can use to retrieve the otherwise
hidden answer session.
\end{tcolorbox}
\captionof{figure}{Case study showing how first-round bridge evidence is transformed into a second-round temporal retrieval cue.}

